\chapter{Häufige Fehler und Stolpersteine}
\label{cha:Stolpersteine}

Dieses Kapitel behandelt die häufigsten Fehler und Stolpersteine der deutschen Grammatik.

\section{Häufige Kasus-Fehler}
\label{sec:KasusFehler}

\begin{enumerate}
    \item \textbf{Seit + Dativ}: Ich lerne seit zwei Jahren. (nicht: Jahr)
    \item \textbf{Anrufen + Akk}: Kannst du mich anrufen? (nicht: mir)
    \item \textbf{glauben + Dativ}: Ich glaube dir. (nicht: dir glauben)
    \item \textbf{helfen + Dativ}: Ich helfe dir. (nicht: dir helfen)
    \item \textbf{antworten + Dativ}: Ich antworte ihm. (nicht: auf ihn antworten - archaish)
\end{enumerate}

\section{Häufige Präpositionen-Fehler}
\label{sec:PraepFehler}

\begin{enumerate}
    \item \textbf{denken an + Akk}: Ich denke an dich. (nicht: dir)
    \item \textbf{warten auf + Akk}: Ich warte auf dich. (nicht: dir)
    \item \textbf{ab + Dativ}: Ich komme ab morgen. (nicht: von morgen)
    \item \textbf{außer + Dativ}: Außer mir war niemand da. (nicht: mich)
\end{enumerate}

\section{Häufige Verb-Fehler}
\label{sec:VerbFehler}

\begin{enumerate}
    \item \textbf{legen/stellen/setzen}:
        \begin{itemize}
            \item Ich lege das Buch auf den Tisch. (horizontal)
            \item Ich stelle das Buch auf den Tisch. (vertikal)
            \item Ich setze mich auf den Stuhl. (mit Bewegung)
        \end{itemize}
    \item \textbf{sein/werden}:
        \begin{itemize}
            \item Er ist müde. (Zustand)
            \item Er wird müde. (Veränderung)
        \end{itemize}
    \item \textbf{haben/sein (Hilfsverben)}:
        \begin{itemize}
            \item Ich habe gegessen. (transitiv)
            \item Ich bin gegangen. (intransitiv/bewegung)
        \end{itemize}
\end{enumerate}

\section{Häufige Artikel-Fehler}
\label{sec:ArtikelFehler}

\begin{enumerate}
    \item \textbf{Nullartikel}: Ich habe Hunger. (nicht: einen Hunger)
    \item \textbf{Nationalitäten}: Er ist Deutscher. (nicht: ein Deutscher - als Substantiv)
    \item \textbf{Berufe}: Sie ist Ärztin. (nicht: eine Ärztin - offiziell)
\end{enumerate}

\section{Häufige Wortstellung-Fehler}
\label{sec:WortstellungFehler}

\begin{enumerate}
    \item \textbf{Deshalb + Inversion}: ..., \textbf{deshalb bin ich} nicht gekommen. (nicht: deshalb ich bin nicht gekommen)
    \item \textbf{Wegen + Inversion}: Wegen des Wetters \textbf{bin ich} geblieben.
    \item \textbf{Relativpronomen}: Das ist der Mann, den ich gesehen habe. (nicht: welcher)
\end{enumerate}

\section{Häufige Zeiten-Fehler}
\label{sec:ZeitenFehler}

\begin{enumerate}
    \item \textbf{Perfekt vs. Präteritum}: In der gesprochenen Sprache wird Perfekt bevorzugt.
    \item \textbf{Futur I vs. Präsens}: Für Zukunft wird oft Präsens verwendet: Morgen fliege ich.
    \item \textbf{Plusquamperfekt}: Nur in Verbindung mit Perfekt/Präteritum: Nachdem ich gegessen hatte, ging ich.
\end{enumerate}

\section{Übersicht: Wichtige Unterscheidungen}
\label{sec:Unterscheidungen}

\begin{center}
\begin{tabular}{|l|l|l|}
\hline
\textbf{A} & \textbf{B} & \textbf{Beispiel} \\ \hline
kennen & wissen & Ich kenne ihn. / Ich weiß, dass... \\ \hline
legen & liegen & Ich lege das Buch. / Das Buch liegt. \\ \hline
stellen & stehen & Ich stelle das Buch. / Das Buch steht. \\ \hline
setzen & sitzen & Ich setze mich. / Ich sitze. \\ \hline
hängen & hängen & Ich hänge das Bild. / Das Bild hängt. \\ \hline
legen & liegen & Ich lege mich ins Bett. / Ich liege im Bett. \\ \hline
bringen & tragen & Ich bringe das Buch. / Ich trage das Hemd. \\ \hline
brauchen & müssen & Ich brauche das nicht. / Ich muss das nicht. \\ \hline
werden & sein & Es wird kalt. / Es ist kalt. \\ \hline
\end{tabular}
\end{center}

\chapterend